% Propietario preservado y actualizado por los cierres sucesores.
\section{Self-inertia as connection endogenous of the joint state}

Let \(\nabla^{\rm tot}\) be a connection on the space of joint states.
The local frame of the observer is obtained by restriction:
\begin{equation}
 \mathfrak F_{\mathcal O}(\Gamma)
 =\operatorname{Res}_{\mathcal O}
 (\Gamma,\nabla^{\rm tot}).
 \label{eq:rm-obs-induced-frame}
\end{equation}
A complete trajectory \(\Gamma(t)\) is self-inertial when
\begin{equation}
 \nabla^{\rm tot}_{\dot\Gamma}\dot\Gamma=0.
 \label{eq:rm-obs-total-geodesic}
\end{equation}
Let \(\pi_S\) be the projection to the observed subsystem and
\(\gamma_S=\pi_S\circ\Gamma\). We define the projection defect
\begin{equation}
 \mathcal K_{\pi_S}(\dot\Gamma,\dot\Gamma)
 :=\nabla^S_{\dot\gamma_S}\dot\gamma_S
 -\dd\pi_S\!\left(
 \nabla^{\rm tot}_{\dot\Gamma}\dot\Gamma
 \right).
 \label{eq:rm-obs-projection-defect}
\end{equation}
On an autoinertial trajectory lies
\begin{equation}
 \boxed{
 a_S=\nabla^S_{\dot\gamma_S}\dot\gamma_S
 =\mathcal K_{\pi_S}(\dot\Gamma,\dot\Gamma).}
 \label{eq:rm-obs-projected-acceleration}
\end{equation}
The acceleration of the subsystem measures the coupling and the degrees of freedom
omitted by the projection. The whole possesses an endogenous connection;
its shadows may exhibit acceleration, rotation, and effective force.

\begin{definition}[Self-inertial state]
TO joint state is self-inertial with respect to a connection and to family
of observers when its evolution satisfies
\cref{eq:rm-obs-total-geodesic} and each local frame is obtained by
\cref{eq:rm-obs-induced-frame}. The apparent forces of each subsystem are
recorded by the defects \(\mathcal K_{\pi_S}\).
\end{definition}

The definition integrates observer and observable within a single dynamic.
A change of frame corresponds to another restriction of the same state and the
same connection. The distinction among affine transport, tetrad rotation,
torsion, and curvature remains available in the geometric codomain.

\section{Machian Criterion of Global Factorization}

Let \(b:\mathcal X_\infty\to B_\infty\) be the global boundary trace and let
\(I_v:\mathcal X_\infty\to V_v\) be a local inertial response. The
response is Machian when there exists
\begin{equation}
 \overline I_v:\operatorname{im}b\longrightarrow V_v,
 \qquad
 I_v=\overline I_v\circ b.
 \label{eq:rm-obs-mach-factorization}
\end{equation}

\begin{theorem}[Machian Response Criterion]
\label{thm:rm-obs-mach-criterion}
There is a unique map \(\overline I_v\) satisfying
\cref{eq:rm-obs-mach-factorization} if and only if \(I_v\) is constant
on each fiber of \(b\):
\begin{equation}
 b(x)=b(x')\quad\Longrightarrow\quad I_v(x)=I_v(x').
 \label{eq:rm-obs-mach-fiber-criterion}
\end{equation}
If \(b\) is surjective, \(\overline I_v\) is defined on all
\(B_\infty\).
\end{theorem}

\begin{proof}
The factorization implies constancy on the fibers. Conversely, for
\(\beta\in\operatorname{im}b\), we define
\(\overline I_v(\beta)=I_v(x)\) with \(x\in b^{-1}(\beta)\). Constancy
guarantees independence of the representative, and the factorization equation
guarantees uniqueness.
\end{proof}

The criterion translates Mach's principle into a verifiable property: the
local response is a reading of the global closure class. To prove
genuinely global dependence, one additionally exhibits two states with the same
local restriction and different boundary traces whose responses differ.
This condition separates an effective global factorization from a response
determined exclusively by the local neighborhood.

The gravitational scale \(Q\) of the CT108 clock provides a dimensionless realization
of this reading:
\begin{equation}
 \alpha_G(E_1,E_2)=\frac{E_1E_2}{Q^2}.
 \label{eq:rm-obs-global-clock-coupling}
\end{equation}
The global clock fixes \(Q\), each local section publishes \(E/Q\), and the
geometry responds through the compliance \(G/c^4\). The equation
provides a Machian chart of the gravitational response; the
\cref{thm:rm-obs-mach-criterion} determines the factorization requirement
that its realization on the complete state must satisfy.

\section{Holonomy and hierarchy of returns}

Let \(\gamma\) be a based loop in \(x_0\), and let
\begin{equation}
 \mathscr H_\gamma=\Hol_\nabla(\gamma)
 \label{eq:rm-obs-holonomy}
\end{equation}
the holonomy in the fiber \(E_{x_0}\). A transported datum returns as
a vector exactly when it belongs to
\begin{equation}
 \operatorname{Fix}(\mathscr H_\gamma)
 =\{\psi:\mathscr H_\gamma\psi=\psi\}.
 \label{eq:rm-obs-fix-holonomy}
\end{equation}
The self-inertial interpretation is condensed in this equality: the frame
transported by the connection of the whole preserves the datum situated in
the fixed subspace.

For a subalgebra of accessible observables
\(\mathcal A_{\mathcal O}\), we define
\begin{align}
 \mathcal R_{\rm vec}
 &=\{\psi:\mathscr H_\gamma\psi=\psi\},
 \label{eq:rm-obs-return-vector}\\
 \mathcal R_{\rm ray}
 &=\{\psi:\mathscr H_\gamma\psi=\ee^{i\theta}\psi
 \text{ for some }\theta\},
 \label{eq:rm-obs-return-ray}\\
 \mathcal R_{\rm obs}
 &=\{\psi:
 \langle\mathscr H_\gamma\psi,A\mathscr H_\gamma\psi\rangle
 =\langle\psi,A\psi\rangle
 \text{ for every }A\in\mathcal A_{\mathcal O}\}.
 \label{eq:rm-obs-return-observable}
\end{align}

\begin{proposition}[Vector--ray--observable hierarchy]
\label{prop:rm-obs-return-hierarchy}
We have
\begin{equation}
 \boxed{
 \mathcal R_{\rm vec}
 \subseteq\mathcal R_{\rm ray}
 \subseteq\mathcal R_{\rm obs}.}
 \label{eq:rm-obs-return-hierarchy}
\end{equation}
\end{proposition}

\begin{proof}
A fixed vector is a fixed ray with zero phase. If
\(\mathscr H_\gamma\psi=\ee^{i\theta}\psi\), the conjugate phases
cancel in every quadratic expectation, so all accessible observables
preserve their value.
\end{proof}

The monodromy nonadic naturally occupies this hierarchy:
\begin{equation}
 g(k+9)=g(k),
 \qquad
 B_{k+9}\neq B_k\quad\hbox{in general}.
 \label{eq:rm-obs-phase-state-return}
\end{equation}
The phase return belongs to the observable reading; the change in
\(B_k\) preserves the memory increment. An injective realization of the
complete state in the fiber makes it possible to lift the criterion
\(\operatorname{Fix}(\mathscr H_\gamma)\) to total genealogical return.
In the absence of that injectivity, the criterion asserts the return of the
realized datum.

\section{Möbius band and conjugation of orientations}

On a dodecaphasic word
\(\mathbf t=(t_0,\ldots,t_{11})\), the genealogical involution is
\begin{equation}
 C_M(\mathbf t)=-\operatorname{rev}(\mathbf t),
 \qquad
 C_M^2=I.
 \label{eq:rm-obs-word-involution}
\end{equation}
Let \(\ell_{\rm ret}\) be an enriched return loop whose lift
has the sign character
\begin{equation}
 \varepsilon_x:\langle\ell_{\rm ret}\rangle\simeq\mathbb Z
 \longrightarrow\{\pm1\},
 \qquad
 \varepsilon_x(\ell_{\rm ret})=-1.
 \label{eq:rm-obs-sign-character}
\end{equation}
The associated interval bundle is
\begin{equation}
 \mathcal M_x
 =([0,1]\times[-1,1])/(0,u)\sim(1,-u).
 \label{eq:rm-obs-mobius-bundle}
\end{equation}

\begin{theorem}[Relative return of Möbius]
\label{thm:rm-obs-mobius-return}
Under the compatible lift of \(C_M\), the loop and the character of
\cref{eq:rm-obs-sign-character}, \(\mathcal M_x\) has nonzero first
Stiefel--Whitney class. One circuit exchanges the local orientations, and
two circuits restore the orientation of the section.
\end{theorem}

\begin{proof}
The transition function of the quotient
\cref{eq:rm-obs-mobius-bundle} is the sign representation of
\(\pi_1(S^1)\). Its class \(w_1\) is the nontrivial generator of
\(H^1(S^1;\mathbb Z/2\mathbb Z)\). The first iteration multiplies the
fiber coordinate by \(-1\) and the second by \((-1)^2=1\).
\end{proof}

The formulation preserves three distinct returns: phase closure
\(C_9\), orientation return in the double cover, and the increment of
memory in the lift \(\mathbb Z\). One turn may restore the
ray and change the vector; two turns restore the orientation and may
preserve a different depth. The band expresses precisely the
difference between publication and genealogy.

The even and odd sectors of a function \(f\) on the double cover are
\begin{equation}
 f^+=\frac12(f+C_M^*f),
 \qquad
 f^-=\frac12(f-C_M^*f).
 \label{eq:rm-obs-even-odd-decomposition}
\end{equation}
The even part descends to the base; the odd part is a section of the
sign line. Quadratic memory resides in the even sector, whereas
orientation and linear charge reside in the odd sector.

\section{Antiunitary conjugation and temporal orientation}

The word involution \(C_M\), the central inversion of the atlas, and the
antiunitary conjugation of the realization in physical space have
distinct domains. Let \(C\) be the antiunitary conjugation on a Hilbert
realization, let \(H=H^*\) be the Hamiltonian, and let \(J_E\) be the complex structure
in the realification. Suppose
\begin{equation}
 CJ_EC^{-1}=-J_E,
 \qquad
 CHC^{-1}=H,
 \qquad
 J_EH\subset HJ_E,
 \label{eq:rm-obs-antiunitary-hypotheses}
\end{equation}
with \(J_E^*=-J_E\), \(J_E^2=-I\), and preservation of the domain of \(H\).
Commutation expresses the complex linearity of the dynamics with respect to
the structure \(J_E\).

\begin{proposition}[Antiunitary reversal of evolution]
\label{prop:rm-obs-antiunitary-time}
For
\begin{equation}
 U(t)=\exp(-J_EHt/\hbar_{\rm int}),
 \label{eq:rm-obs-unitary-evolution}
\end{equation}
we have
\begin{equation}
 \boxed{CU(t)C^{-1}=U(-t).}
 \label{eq:rm-obs-time-reversal}
\end{equation}
\end{proposition}

\begin{proof}
The generator \(-J_EH/\hbar_{\rm int}\) is anti-self-adjoint. The hypotheses of
\cref{eq:rm-obs-antiunitary-hypotheses} change its sign under conjugation.
Functional calculus and uniqueness of the generated unitary group yield
\cref{eq:rm-obs-time-reversal}.
\end{proof}

When, additionally, a charge operator \(Q_{\rm ch}\) and a mass
observable \(M\) satisfy
\begin{equation}
 CQ_{\rm ch}C^{-1}=-Q_{\rm ch},
 \qquad
 CMC^{-1}=M,
 \label{eq:rm-obs-charge-mass-parity}
\end{equation}
the two orientations of the band are published as carriers of conjugate charge
and shared even memory. The particle--antiparticle representation requires the intertwiner that lifts particle--
antiparticle \(C_M\) to \(C\) and
preserves the evolution. The identities
\cref{eq:rm-obs-word-involution,eq:rm-obs-time-reversal} provide the
two typed endpoints of that intertwiner.

The tritic temporal reader organizes the regimes according to
\begin{equation}
 \begin{aligned}
 +1&\longmapsto\text{return, memory and past},\\
 0&\longmapsto\text{torsional threshold and present},\\
 -1&\longmapsto\text{bidirectional opening and future}.
 \end{aligned}
 \label{eq:rm-obs-temporal-trit}
\end{equation}
The assignment is to reader of the structure algebraic
\(J^2=-1,0,+1\). The promotion to curvatures of a spacetime requires a
connection and a chart metric; the temporal reader meanwhile preserves
its exact meaning of return, threshold, and opening.

\section{Recoverability and Landauer's Principle}

Landauer's thermodynamic principle locates the physical cost of
logically irreversible erasure \cite{Landauer1961}.

Let \((\Omega,2^\Omega,\mu)\) to finite probabilized space, let
\(X:\Omega\to\mathcal X\) the complete state variable and let
\(q:\mathcal X\to Y\) to reader. The chain rule provides
\begin{equation}
 \mathsf H(X)
 =\mathsf H(q(X))+\mathsf H(X\mid q(X)).
 \label{eq:rm-obs-entropy-chain}
\end{equation}
The second summand is the information contained in the reader fiber.

\begin{theorem}[Landauer zero of full transport]
\label{thm:rm-obs-zero-landauer}
If \(U:\mathcal X\to\mathcal X\) is a bijective update, then
\begin{equation}
 \boxed{\mathsf H(X\mid U(X))=0.}
 \label{eq:rm-obs-zero-conditional-entropy}
\end{equation}
The logical information erased by the update is null. For every
physical protocol \(\mathcal P_U\) that realizes it, the Landauer bound and its
ideal reversible limit satisfy
\begin{equation}
 \boxed{
 Q_{\rm L}[\mathcal P_U]\geq0,
 \qquad
 \inf_{\mathcal P_U^{\rm rev}}Q_{\rm L}[\mathcal P_U]=0.}
 \label{eq:rm-obs-zero-landauer-cost}
\end{equation}
\end{theorem}

\begin{proof}
The inverse of \(U\) reconstructs \(X\) from \(U(X)\); therefore, the conditional distribution is concentrated on a single state. Its entropy is zero, and the operation, which is logical, preserves information. The resulting thermodynamic bound is \(Q_{\rm L}\geq0\); reversible quasistatic protocols attain the ideal null infimum. This statement distinguishes the bound from the concrete dissipation of a device.
\end{proof}

The projected output \(q(U(X))\) preserves part of the information and
leaves in the fiber
\begin{equation}
 \mathsf H(X\mid q(U(X))).
 \label{eq:rm-obs-reader-forgotten-info}
\end{equation}
Let \(M:\mathcal X\to\mathfrak M\) be a memory register, and define the
enriched reader
\begin{equation}
 \widehat q(x)=(q(x),M(x)).
 \label{eq:rm-obs-enriched-reader}
\end{equation}

\begin{proposition}[Recoverability via output and memory]
\label{prop:rm-obs-recoverability}
If \(\widehat q\) is injective on the support of \(X\), then
\begin{equation}
 \boxed{\mathsf H(X\mid q(X),M(X))=0.}
 \label{eq:rm-obs-recoverability}
\end{equation}
The quantity \(\mathsf H(X\mid q(X))\) measures the additional memory required to
recover the state from the visible output.
\end{proposition}

\begin{proof}
Injectivity on the support provides a deterministic inverse from
\(\widehat q(X)\) to \(X\). The conditional entropy is therefore zero. The
chain rule locates the difference between the visible reader and the
enriched reader in the fiber of \(q\).
\end{proof}

For the reset of a uniform trit, the logical entropy is \(\log3\),
and the Landauer bound takes the form
\begin{equation}
 \boxed{Q_{\rm reset}\geq k_B\Theta\log3.}
 \label{eq:rm-obs-trit-reset}
\end{equation}
The structural identity of the clock,
\begin{equation}
 E_P=k_B\Theta_{\rm clk}\log3,
 \label{eq:rm-obs-planck-trit-landauer}
\end{equation}
evaluates that bound at the declared chronological temperature. The Planck
quantum then appears as the minimum cost of resetting a uniform ternary decision
in that thermal chart.

The zero Landauer cost and the cost \(k_B\Theta\log3\) describe different
operations. The former corresponds to the bijective transport of the
enriched state; the latter corresponds to a reset that identifies three
registers. The costs of control, finite speed, preparation, and
coupling to the environment belong to the physical realization of the
device. The logical equality of zero minimum cost preserves
information; an equation of motion for a bounce or a
cosmological cycle additionally requires boundary dynamics and its balance laws.

\section{Reversible measurement, conditioning, and reinitialization}

The operations of the complete process are ordered as
\begin{equation}
 \boxed{\begin{gathered}
 \text{preparation}
 \longrightarrow
 \text{reversible coupling}
 \longrightarrow
 \text{pointer reading}\\
 \longrightarrow
 \text{fiber conditioning}
 \longrightarrow
 \text{reinitialization}
 \end{gathered}}
 \label{eq:rm-obs-measurement-chain}
\end{equation}
Each arrow has its own informational balance:
\begin{enumerate}[label=\textup{(\roman*)}]
 \item Unitary premeasurement preserves purity and norm
in the composite system.
 \item The reading publishes \(q(X)\) and preserves the record \(M(X)\) in the
 apparatus.
 \item Conditioning selects the relative fiber and normalizes its
 measure.
 \item The nonselective channel sums all instruments and preserves the
 trace.
 \item Resetting identifies registers and activates the Landauer bound
 corresponding to the erased information.
\end{enumerate}

This chain localizes the origin of operational irreversibility. The global
state may evolve isometrically while a partial reading
exhibits increasing entropy. Information accessible to one branch and
information preserved in global correlations are distinct
magnitudes. The coupled observer forms part of those correlations, and its
memory constitutes a physical component of the joint state.

\section{Confluence of observation, self-inertia, and modular response}

The above results are assembled into the causal chain.
\begin{equation}
 \begin{aligned}
 \text{enriched APP--TRIT--TPK state}
 &\longrightarrow
 (\text{section},\text{observer},\text{memory})\\
 &\longrightarrow
 (\text{fiber},\text{PVM},\text{instrument})\\
 &\longrightarrow
 (\text{relative state},\text{holonomy},\text{induced frame})\\
 &\longrightarrow
 (\text{Machian response},\text{Typed Landauer}).
 \end{aligned}
 \label{eq:rm-obs-causal-chain}
\end{equation}
The modular structure of the preceding chapter adds the pair
\begin{equation}
 (\widehat{\mathcal A},K_A),
 \label{eq:rm-obs-area-modular-pair}
\end{equation}
which preserves geometric quantity and provenance. The observer accesses a
compression of this pair through its PVM; the memory preserves the channel from
which the result originates; autoinertia transports the frame along the
connection; relative entropy measures the cost of changing genealogical
state; and Landauer localizes the cost of the effective erasure of the
record.

The relation to Everett lies at the first level of this chain: the isometry preserves the superposition of registers, and the observer occupies a relative state. The relation to Mach lies at the connection level: the frame and the local response are induced from the compatible totality. The Möbius band lies at the holonomy level: the orientation can change while the ray and the accessible observables return. The three readings share the enriched state and preserve different mathematical functions.

\section{Domain of validity and physical realization conditions}

The probative force is distributed as follows:
\begin{enumerate}[label=\textup{(\roman*)}]
 \item The fibre--PVM--Lüders correspondence, the Kraus representation,
 the CPTP channel, and the dilation are
 \status{exact results retrieved} in the declared domains.
 \item The formulation of the relative state and the conservation, isometric in character,
 of the channel are \status{exact sciences}. The Everettian ontology of branches and their
 macroscopic stability belong to a
 \status{interpretive and dynamic interface}.
 \item The criterion
 \(\operatorname{Fix}(\Hol)\), the vector--ray--observable hierarchy
 and the Machian factorization are
 \status{exact results retrieved}. Its application to a gravitational geometry
 requires the corresponding realization morphism.
 \item The Möbius band is \status{exacta relativamente} to the
 lift, to the loop, and to the declared sign character. The
 interpretation of particle--antiparticle requires the antiunitary intertwiner
 of \cref{eq:rm-obs-antiunitary-hypotheses}.
 \item The zero Landauer cost of a bijective update, the recoverability
 of the enriched reader, and the ternary reset bound are
 \status{exact results retrieved}. The cosmological selector and the
 dynamics of sheets prolong these balances by means of the declared equations of
 evolution, splicing conditions, and torsional memory.
\end{enumerate}

\begin{proofstatusbox}
The chapter proves in its finite domains the equivalence of conditionings, preservation of the global channel, the return and factorization criteria, and the logical Landauer balances. The joint realization employs the isometric intertwiner between the TPK memory and the apparatus, the dynamic stability of the finite register, and transport of the genealogical involution to antiunitary conjugation with its even and odd observables.
\end{proofstatusbox}

\begin{falsifierbox}
The corresponding statements are refuted by: a nonmeasurable fiber; a family
of projectors lacking a common PVM; a Lüders conditioning
that differs from \cref{eq:rm-obs-luders-fiber-equivalence}; Kraus
operators whose sum of effects differs from the identity; a dilation with
\(V_a^\dagger V_a\neq I\); an exact vector return attributed to a datum
outside \(\operatorname{Fix}(\Hol)\); a Machian response varying within a
fiber of \(b\); a loop with a trivial sign character presented as a Möbius
strip; a conjugation that fails to satisfy
\(CU(t)C^{-1}=U(-t)\); a bijective update with positive conditional
entropy; or a uniform ternary reset whose balance omits the
information \(\log3\). A cosmological promotion is also refuted by
presenting zero Landauer as a dynamical bounce equation.
\end{falsifierbox}
